\clearpage
\begin{proofbox}
\textbf{统一记号：}沿用结论中的双曲线及焦点，记 $r_1=PF_1>0$、$r_2=PF_2>0$，
则 $|r_1-r_2|=2a$，$F_1F_2=2c$。

\medskip
\textbf{方法一：最快证明【定义＋三角形不等式，2--3 步】。}
\begin{center}\begin{tikzpicture}[font=\small,>=stealth]
\draw[gray!45,rounded corners=3pt] (0,0) rectangle (14.2,10.5);
\node at (7.1,9.85) {方法一：用同一单位，把两腰之和、差与底边作比较};

\begin{scope}[shift={(7.1,6.4)},scale=.95]
\fill[orange!7] (-2,0)--(1.5,{sqrt(15)/2})--(2,0)--cycle;
\draw[orange!85!black,very thick] (-2,0)--(1.5,{sqrt(15)/2});
\draw[purple!80!black,very thick] (1.5,{sqrt(15)/2})--(2,0);
\draw[teal!75!black,thick] (-2,0)--(2,0);
\fill (-2,0) circle(3pt) node[below left=5pt] {$F_1$};
\fill (2,0) circle(3pt) node[below right=5pt] {$F_2$};
\fill (1.5,{sqrt(15)/2}) circle(3pt) node[above=5pt] {$P$};
\node[orange!85!black,left=5pt] at (-.25,{sqrt(15)/4}) {$r_1=4$};
\node[purple!80!black,right=5pt] at (1.75,{sqrt(15)/4}) {$r_2=2$};
\node[teal!75!black,below=5pt] at (0,0) {$2c=4$};
\end{scope}
\draw[orange!85!black,line width=3pt] (3.2,4.45)--(7,4.45);
\draw[purple!80!black,line width=3pt] (7,4.45)--(8.9,4.45);
\node[anchor=east] at (2.8,4.45) {和};\node[anchor=west] at (9.3,4.45) {$6\ge2c=4$};
\draw[orange!85!black,line width=3pt] (3.2,3.2)--(5.1,3.2);
\node[anchor=east] at (2.8,3.2) {差};\node[anchor=west] at (5.5,3.2) {$2=2a$};
\node at (7.1,1.75) {$\displaystyle e=\frac{2c}{2a}\le\frac{r_1+r_2}{|r_1-r_2|}$};
\node at (7.1,.65) {图取 $a=1,c=2,\lambda=2$；上方三角形与下方线段使用同一长度单位。};

\end{tikzpicture}
\end{center}
\textbf{先认图：} 两腰的差固定为 $2a$，两腰的和至少覆盖底边 $2c$。

\textbf{第1步：}把双曲线定义与三角形条件并排写出：
\[
\boxed{|r_1-r_2|=2a},\qquad
\boxed{r_1+r_2\ge2c}.
\]
\textbf{第2步：}两式作比（分母为正），并直接代入 $r_1=\lambda r_2$：
\[
\frac ca
\le\frac{r_1+r_2}{|r_1-r_2|}
=\frac{\lambda r_2+r_2}{|\lambda r_2-r_2|}
=\boxed{\frac{1+\lambda}{|1-\lambda|}}.
\]
\textbf{第3步：}由 $e>1$ 得
\[
\boxed{1<e\le\frac{1+\lambda}{|1-\lambda|}},\qquad \lambda\ne1.
\]
若 $\lambda=1$，则 $r_1=r_2$ 与 $|r_1-r_2|=2a>0$ 矛盾。

\textbf{等号及可达性：}等号时 $P$ 在两焦点间的对应顶点：$\lambda>1$ 取右顶点，
$0<\lambda<1$ 取左顶点。范围内确有点由方法五的两圆构造保证，故为\textbf{充要条件}。
\textbf{最快的原因：}不设坐标，只把两腰的和与差直接作比。
\end{proofbox}

\clearpage
\begin{proofbox}
\textbf{方法二：最低门槛证明【初中分份法：差定求和】。}
\begin{center}\begin{tikzpicture}[font=\small,>=stealth]
\draw[gray!45,rounded corners=3pt] (0,0) rectangle (14.2,9.5);
\node at (7.1,8.85) {方法二：两腰按份数画出，差定就能求和};

\node at (7.1,7.9) {先取 $\lambda=2$：长腰两份、短腰一份。};
\draw[orange!85!black,line width=4pt] (1.1,6.7)--(9.1,6.7);
\foreach \x in {1.1,5.1,9.1}\draw (\x,6.5)--(\x,6.9);
\node[anchor=east] at (.9,6.7) {$r_1$};
\draw[purple!80!black,line width=4pt] (1.1,4.6)--(5.1,4.6);
\foreach \x in {1.1,5.1}\draw (\x,4.4)--(\x,4.8);
\node[anchor=east] at (.9,4.6) {$r_2$};
\draw[<->] (5.1,5.9)--(9.1,5.9) node[midway,below=4pt] {差：$1$ 份 $=2a$};
\node[anchor=west] at (6.1,4.6) {所以一份长 $2a$};
\draw[orange!85!black,line width=4pt] (1.1,2.5)--(9.1,2.5);
\draw[purple!80!black,line width=4pt] (9.1,2.5)--(13.1,2.5);
\node at (7.1,1.7) {和：$3$ 份 $=6a$，必须不小于 $2c$};
\node at (7.1,.65) {一般情形：差 $|\lambda-1|$ 份、和 $\lambda+1$ 份。};

\end{tikzpicture}
\end{center}
\textbf{先从整数份看：} 图中 $2:1$ 的两腰，差一份、和三份；再推广到 $\lambda:1$。

\textbf{第1步：}已知 $r_1:r_2=\lambda:1$，就设
\[
r_1=\lambda t,\qquad r_2=t,\qquad t>0.
\]
这只引入\textbf{一个未知数} $t$；无需点坐标、参数方程或任何高等技巧。

\textbf{第2步：}双曲线规定\textbf{两腰之差}为 $2a$：
先排除 $\lambda=1$，因为两腰相等会与 $2a>0$ 矛盾。以下 $\lambda\ne1$。
\[
|\lambda t-t|=2a
\quad\Longrightarrow\quad
\boxed{t=\frac{2a}{|\lambda-1|}}.
\]
例如 $\lambda=2$ 时，腰长比 $2:1$，两腰之差就是“$1$ 份”$=2a$；
因此两腰之和是“$3$ 份”$=6a$。

\textbf{第3步：}底边是两焦点间距 $2c$，必须满足
\[
2c\le r_1+r_2=(\lambda+1)t
=\frac{2a(\lambda+1)}{|\lambda-1|}.
\]
除以 $2a$，得到
\[
\boxed{1<e=\frac ca\le\frac{\lambda+1}{|\lambda-1|}}.
\]
\textbf{何时等号：}两腰首尾拼成底边，即 $P$ 落在焦点之间的顶点；
否则 $PF_1+PF_2>F_1F_2$，上界取不到。

\textbf{最低门槛在哪里：}全程只需\textbf{比例、绝对值、三角形两边之和}；
考场忘了公式也可从最熟悉的“差几份、和几份”当场还原。
\end{proofbox}

\clearpage
\begin{proofbox}
\textbf{方法三：常规解析几何【焦半径公式＋坐标范围】。}
\begin{center}\begin{tikzpicture}[font=\small,>=stealth]
\draw[gray!45,rounded corners=3pt] (0,0) rectangle (14.2,11.5);
\node at (7.1,10.85) {方法三：先判左右支，再看竖线能否碰到曲线};

\node at (3.6,9.75) {$\lambda=1/2$：左支};\node at (10.6,9.75) {$\lambda=2$：右支};
\begin{scope}[shift={(7.1,5.6)},scale=.8]
\draw[gray,->] (-5.4,0)--(5.4,0) node[right] {$x$};
\draw[gray,->] (0,-4.7)--(0,5) node[above] {$y$};

\draw[teal!75!black,thick,samples=90,domain=-4.2:4.2,variable=\v] plot ({2*sqrt(1+\v*\v/12)},{\v});
\draw[teal!75!black,thick,samples=90,domain=-4.2:4.2,variable=\v] plot ({-2*sqrt(1+\v*\v/12)},{\v});

\draw[purple!80!black,dashed,thick] (-3,-4.5)--(-3,4.5);
\draw[orange!85!black,dashed,thick] (3,-4.5)--(3,4.5);
\foreach \sgn in {-1,1}{\fill[teal!75!black] ({3*\sgn},{sqrt(15)}) circle(3pt);\fill[teal!75!black] ({3*\sgn},{-sqrt(15)}) circle(3pt);}
\node[above right=5pt] at (3,{sqrt(15)}) {$P$};\node[above left=5pt] at (-3,{sqrt(15)}) {$Q$};
\fill (-4,0) circle(2.8pt) node[below left=5pt] {$F_1$};\fill (4,0) circle(2.8pt) node[below right=5pt] {$F_2$};
\fill (-2,0) circle(2.8pt) node[below right=5pt] {$-a$};\fill (2,0) circle(2.8pt) node[below left=5pt] {$a$};
\node[purple!80!black,anchor=east] at (-3.2,.8) {$x=-3$};
\node[orange!85!black,anchor=west] at (3.2,.8) {$x=3$};
\draw[gray,dotted,<->] (-2.6,{sqrt(15)})--(2.6,{sqrt(15)});
\node[above=4pt,fill=white,inner sep=2pt] at (0,{sqrt(15)}) {关于 $y$ 轴反射};
\end{scope}
\node at (7.1,1.25) {右支要求 $x\ge a$；左支要求 $x\le-a$；等号时竖线只碰到顶点。};
\node at (7.1,.6) {本图 $a=2,c=4$；$x=\pm3$ 均给出 $y=\pm\sqrt{15}$。};

\end{tikzpicture}
\end{center}
\textbf{先看图中的竖线：} 横坐标落在相应支的取值范围内，才有真实交点。

在\textbf{右支}上 $x\ge a$，两焦半径为
\[
PF_1=ex+a,\qquad PF_2=ex-a.
\]
故 $\lambda=(ex+a)/(ex-a)>1$。解得
\[
\boxed{x=\frac ae\,\frac{\lambda+1}{\lambda-1}}.
\]
而双曲线右支存在点 $P$ 的充要条件正是 $x\ge a$；
因此
\[
\lambda>1:\quad
\frac{\lambda+1}{\lambda-1}\ge e>1.
\]

\tcbbreak
\textbf{方法三续：左支是右支关于 $y$ 轴的镜像。}
\begin{center}\begin{tikzpicture}[font=\small,>=stealth]
\draw[gray!45,rounded corners=3pt] (0,0) rectangle (14.2,11.5);
\node at (7.1,10.85) {方法三：先判左右支，再看竖线能否碰到曲线};

\node at (3.6,9.75) {$\lambda=1/2$：左支};\node at (10.6,9.75) {$\lambda=2$：右支};
\begin{scope}[shift={(7.1,5.6)},scale=.8]
\draw[gray,->] (-5.4,0)--(5.4,0) node[right] {$x$};
\draw[gray,->] (0,-4.7)--(0,5) node[above] {$y$};

\draw[teal!75!black,thick,samples=90,domain=-4.2:4.2,variable=\v] plot ({2*sqrt(1+\v*\v/12)},{\v});
\draw[teal!75!black,thick,samples=90,domain=-4.2:4.2,variable=\v] plot ({-2*sqrt(1+\v*\v/12)},{\v});

\draw[purple!80!black,dashed,thick] (-3,-4.5)--(-3,4.5);
\draw[orange!85!black,dashed,thick] (3,-4.5)--(3,4.5);
\foreach \sgn in {-1,1}{\fill[teal!75!black] ({3*\sgn},{sqrt(15)}) circle(3pt);\fill[teal!75!black] ({3*\sgn},{-sqrt(15)}) circle(3pt);}
\node[above right=5pt] at (3,{sqrt(15)}) {$P$};\node[above left=5pt] at (-3,{sqrt(15)}) {$Q$};
\fill (-4,0) circle(2.8pt) node[below left=5pt] {$F_1$};\fill (4,0) circle(2.8pt) node[below right=5pt] {$F_2$};
\fill (-2,0) circle(2.8pt) node[below right=5pt] {$-a$};\fill (2,0) circle(2.8pt) node[below left=5pt] {$a$};
\node[purple!80!black,anchor=east] at (-3.2,.8) {$x=-3$};
\node[orange!85!black,anchor=west] at (3.2,.8) {$x=3$};
\draw[gray,dotted,<->] (-2.6,{sqrt(15)})--(2.6,{sqrt(15)});
\node[above=4pt,fill=white,inner sep=2pt] at (0,{sqrt(15)}) {关于 $y$ 轴反射};
\end{scope}
\node at (7.1,1.25) {右支要求 $x\ge a$；左支要求 $x\le-a$；等号时竖线只碰到顶点。};
\node at (7.1,.6) {本图 $a=2,c=4$；$x=\pm3$ 均给出 $y=\pm\sqrt{15}$。};

\end{tikzpicture}
\end{center}
在\textbf{左支}上写 $x=-t$（$t\ge a$），有
\[
PF_1=et-a,\qquad PF_2=et+a,
\qquad \lambda=\frac{et-a}{et+a}<1.
\]
解得
\[
\boxed{t=\frac ae\,\frac{1+\lambda}{1-\lambda}}.
\]
同理，左支存在点的充要条件为 $t\ge a$，所以
\[
0<\lambda<1:\quad
\frac{1+\lambda}{1-\lambda}\ge e>1.
\]
合并两支即
\[
\boxed{\lambda\ne1,\quad 1<e\le\frac{1+\lambda}{|1-\lambda|}}.
\]
\textbf{充分性不靠猜图：}以上条件成立时，右支的 $x\ge a$ 或左支的
$t\ge a$，代入
\[
\boxed{y=\pm\frac ba\sqrt{x^2-a^2}}
\]
得到真实曲线点；当 $|x|>a$ 有\textbf{两个}对称解，$|x|=a$ 只有顶点一解。
\textbf{常规法优点：}不仅求范围，还能求坐标、判断所在支及交点个数。
\end{proofbox}

\clearpage
\begin{proofbox}
\textbf{方法四：阿波罗尼斯圆证明【一条定比轨迹圆的长度最值】。}

\textbf{补充理由：}这是\textbf{独立的几何证明}。
方法一只看某一点的三角形，本法先找出\emph{所有满足距离比的点}构成的圆，
再问它能否和双曲线的\emph{定差轨迹}相交。

\textbf{第1步：建立阿波罗尼斯圆。}先设 $\lambda>1$（$0<\lambda<1$ 稍后由对称处理），
令 $P=(x,y)$，则
\[
PF_1=\lambda PF_2
\iff (x+c)^2+y^2=\lambda^2[(x-c)^2+y^2].
\]
整理并配方：
\[
\boxed{(x-x_C)^2+y^2=R^2},\qquad
x_C=\frac{c(\lambda^2+1)}{\lambda^2-1},\quad
R=\frac{2c\lambda}{\lambda^2-1}.
\]
故定比轨迹确是一条阿波罗尼斯圆 $\mathcal A_\lambda$。
\begin{center}% C086 阿波罗尼斯圆与双曲线精确示意（数学坐标，非自由重绘）
% H:x^2/4-y^2/12=1, a=2, c=4, e=2; lambda=2.
% F1=(-4,0), F2=(4,0).
% Apollonius circle center C=(20/3,0), radius R=16/3.
% H cap A=(2,0); circle leftmost T=(4/3,0).
% Exact intersections P+/-=(3,+/-sqrt(15)); PF1=8, PF2=4.
\begin{tikzpicture}[x=0.58cm,y=0.58cm,>=stealth]
  \draw[->,gray!65] (-4.85,0)--(12.55,0) node[right] {$x$};
  \draw[->,gray!65] (0,-5.8)--(0,5.8) node[above] {$y$};
  % Circle locus. The path is constructed from exact center and radius.
  \draw[orange!85!black,very thick] ({20/3},0) circle [radius={16/3}];
  % Entire visible hyperbola branches from analytic formula y=+/-sqrt(3(x^2-4)).
  \draw[teal!80!black,very thick,samples=95,domain=2:3.75,smooth]
     plot (\x,{sqrt(3*(\x*\x-4))});
  \draw[teal!80!black,very thick,samples=95,domain=2:3.75,smooth]
     plot (\x,{-sqrt(3*(\x*\x-4))});
  \draw[teal!80!black,very thick,samples=95,domain=-3.75:-2,smooth]
     plot (\x,{sqrt(3*(\x*\x-4))});
  \draw[teal!80!black,very thick,samples=95,domain=-3.75:-2,smooth]
     plot (\x,{-sqrt(3*(\x*\x-4))});
  \coordinate (Fone) at (-4,0);
  \coordinate (Ftwo) at (4,0);
  \coordinate (Ap) at (2,0);
  \coordinate (Cc) at ({20/3},0);
  \coordinate (T) at ({4/3},0);
  \coordinate (Pup) at (3,{sqrt(15)});
  \coordinate (Pdn) at (3,{-sqrt(15)});
  \draw[dashed,purple!75!black,thin] (Fone)--(Pup)--(Ftwo);
  \fill[black] (Fone) circle (2pt) node[below left] {$F_1$};
  \fill[black] (Ftwo) circle (2pt) node[below=5pt] {$F_2$};
  \fill[teal!75!black] (Ap) circle (2.3pt) node[below right=10pt] {$A$};
  \fill[orange!85!black] (Cc) circle (2pt) node[below=5pt] {$C_\lambda$};
  \fill[orange!85!black] (T) circle (2.2pt) node[above left=8pt] {$T$};
  \fill[red!80!black] (Pup) circle (2.8pt) node[above left] {$P_+$};
  \fill[red!80!black] (Pdn) circle (2.8pt) node[below left] {$P_-$};
  \node[fill=white,fill opacity=0.9,text opacity=1,text=orange!80!black]
    at (8.7,5.65) {$\mathcal A_2$: 定比轨迹圆};
  \node[fill=white,fill opacity=0.9,text opacity=1,text=teal!75!black]
    at (-1.3,5.35) {$H$: 定差轨迹};
  \node[fill=white,fill opacity=0.9,text opacity=1] at (8.6,-4.75)
    {$P_\pm=(3,\pm\sqrt{15})$};
\end{tikzpicture}
\end{center}
\textbf{图意：} 橙色圆给出全部定比点，绿色双曲线要求距离差为 $2a$；公共点才是所求点。

\textbf{交点数和相切位置：}当 $2a$ 严格处于下面将求出的距离差最小值与最大值之间，
圆上恰有上下对称\textbf{两点}满足；等于最小值时只有\textbf{一个}，
且为右顶点 $A=(a,0)$（$0<\lambda<1$ 时为左顶点），两曲线在顶点处\textbf{相切}。
$\lambda=1$ 时轨迹是 $x=0$，与双曲线没有交点。
\textbf{识别记忆：}\emph{定比成圆，定差成双曲线；看两个轨迹能否相交。}

\tcbbreak
\textbf{第2步：在这条圆上求 $PF_2$ 的最小、最大值。}
\begin{center}% lambda=2, c=4: center (20/3,0), R=16/3; scale=.525 cm/unit.
\begin{tikzpicture}[font=\small,>=stealth]
\draw[gray!45,rounded corners=3pt] (0,0) rectangle (14.2,10.4);
\node at (7.1,9.75) {方法四：圆内定点到圆周，最近为 $R-d$，最远为 $R+d$};
\node at (7.1,8.95) {例：$\lambda=2,c=4$，$R=16/3$，$d=CF_2=8/3<R$};
\begin{scope}[shift={(4,5.5)},scale=.525]
\coordinate (C) at (6.6666666667,0);
\coordinate (Ftwo) at (4,0);
\coordinate (T) at (1.3333333333,0);
\coordinate (U) at (12,0);
\draw[orange!85!black,very thick] (C) circle (5.3333333333);
\draw[gray] (.6,0)--(12.7,0);
\draw[purple!80!black,line width=2pt] (T)--(Ftwo);
\draw[teal!75!black,thick] (Ftwo)--(U);
\draw[orange!85!black,dashed,thick] (C)--(6.6666666667,5.3333333333)
 node[midway,right=6pt] {$R$};
\fill (C) circle (2.5pt) node[below=7pt] {$C$};
\fill (Ftwo) circle (2.5pt) node[below=7pt] {$F_2$};
\fill[purple!80!black] (T) circle (3pt) node[above left=6pt] {$T$};
\fill[teal!75!black] (U) circle (3pt) node[above right=6pt] {$U$};
\node at (5.3333333333,.75) {$d$};
\end{scope}
\draw[purple!80!black,<->,thick] (4.7,1.8)--(6.1,1.8)
 node[midway,above=5pt] {$F_2T=R-d=8/3$};
\draw[teal!75!black,<->,thick] (6.1,.7)--(10.3,.7)
 node[midway,above=5pt] {$F_2U=R+d=8$};
\end{tikzpicture}
\end{center}
因 $x_C>c$，右焦点 $F_2(c,0)$ 到圆心的距离为
\[
d=x_C-c=\frac{2c}{\lambda^2-1},\qquad R>d,
\]
所以 $F_2$ 在圆内。
从圆内定点到圆周的距离范围是 $[R-d,R+d]$，因此
\[
\boxed{\frac{2c}{\lambda+1}\le PF_2\le\frac{2c}{\lambda-1}}.
\]
在此轨迹上 $PF_1-PF_2=(\lambda-1)PF_2$，随之得到
\[
\boxed{\frac{2c(\lambda-1)}{\lambda+1}
\le PF_1-PF_2\le2c}.
\]

\textbf{第3步：用双曲线的\emph{定差} $2a$ 与该范围相交。}
由于 $c>a>0$，自动有 $2a<2c$。因此圆上存在点满足
$PF_1-PF_2=2a$ 当且仅当
\[
2a\ge\frac{2c(\lambda-1)}{\lambda+1}
\iff
\boxed{1<e\le\frac{\lambda+1}{\lambda-1}}.
\]
当 $0<\lambda<1$，把 $F_1,F_2$ 互换并使用 $1/\lambda>1$，
得到统一结论
\[
\boxed{\exists P\in H:\ PF_1=\lambda PF_2
\iff \lambda\ne1,\quad1<e\le\frac{1+\lambda}{|1-\lambda|}}.
\]

\end{proofbox}

\clearpage
\begin{proofbox}
\textbf{方法五：几何补充【两焦点辅助圆相交，构造充分性】。}

\textbf{构造目的：}以两个焦点分别作圆心、固定两条半径，直接找满足条件的点；
这与方法四的\textbf{一条定比轨迹圆}不同。

假定 $\lambda>0$ 且 $\lambda\ne1$，并满足
$1<e\le(1+\lambda)/|1-\lambda|$。
\begin{center}\begin{tikzpicture}[font=\small,>=stealth]
\draw[gray!45,rounded corners=3pt] (0,0) rectangle (14.2,8.5);
\node at (7.1,7.85) {方法五：两焦点作圆心，交点个数决定是否有解};

\node at (7.1,6.9) {取 $a=1,\lambda=2$，于是 $r_1=4,r_2=2$；只改变 $c=e$。};
\node at (2.5,6.2) {$e=2$};\begin{scope}[shift={(2.75,4.4)},scale=.3]
\draw[orange!85!black,thick] (-2,0) circle(4);\draw[purple!80!black,thick] (2,0) circle(2);\draw[gray] (-2,0)--(2,0);\fill (-2,0) circle(3pt) node[below=43pt] {$F_1$};\fill (2,0) circle(3pt) node[below=43pt] {$F_2$};
\fill[teal!75!black] (1.5,1.9364916731037085) circle(4pt);\fill[teal!75!black] (1.5,-1.9364916731037085) circle(4pt);\end{scope}
\node at (2.5,1.65) {相交：两个点};
\node at (7.0,6.2) {$e=3$};\begin{scope}[shift={(7.25,4.4)},scale=.3]
\draw[orange!85!black,thick] (-3,0) circle(4);\draw[purple!80!black,thick] (3,0) circle(2);\draw[gray] (-3,0)--(3,0);\fill (-3,0) circle(3pt) node[below=43pt] {$F_1$};\fill (3,0) circle(3pt) node[below=43pt] {$F_2$};
\fill[teal!75!black] (1.0,0.0) circle(4pt);\end{scope}
\node at (7.0,1.65) {外切：一个顶点};
\node at (11.5,6.2) {$e=3.5$};\begin{scope}[shift={(11.75,4.4)},scale=.3]
\draw[orange!85!black,thick] (-3.5,0) circle(4);\draw[purple!80!black,thick] (3.5,0) circle(2);\draw[gray] (-3.5,0)--(3.5,0);\fill (-3.5,0) circle(3pt) node[below=43pt] {$F_1$};\fill (3.5,0) circle(3pt) node[below=43pt] {$F_2$};
\end{scope}
\node at (11.5,1.65) {分离：没有点};

\node at (7.1,.65) {三图同一尺度；圆心距 $2c$ 与半径和 $r_1+r_2=6$ 比较。};

\end{tikzpicture}
\end{center}
\textbf{先看同尺度的三种状态：} 两圆交点就是满足两条指定焦半径的点。一般地取
\[
r_2=\frac{2a}{|\lambda-1|},\qquad
r_1=\frac{2a\lambda}{|\lambda-1|}.
\]
取以 $F_1,F_2$ 为圆心、分别以 $r_1,r_2$ 为半径的两个圆。
两圆圆心距是 $F_1F_2=2c$，而
\[
|r_1-r_2|=2a<2c,\qquad
r_1+r_2=\frac{2a(1+\lambda)}{|1-\lambda|}\ge2c.
\]
这正是两圆有交点的\textbf{充要条件}。任一交点 $P$ 满足
$PF_1=r_1$、$PF_2=r_2$，故
\[
|PF_1-PF_2|=2a,\qquad PF_1=\lambda PF_2,
\]
即 $P$ 在所给双曲线上。
\textbf{对应图中三列：}$r_1+r_2>2c$ 时相交于两个对称点；
$r_1+r_2=2c$ 时外切于相应顶点；$r_1+r_2<2c$ 时分离、无解。
\textbf{特别注意：}两圆相切的正确类型是\textbf{外切}，不是内切；
这是因为 $r_1+r_2=2c$，而 $|r_1-r_2|=2a<2c$。

\end{proofbox}
